\section{Data Recipe: la receta de datos es la verdadera barrera}

Con la misma fuente de datos, distintas formas de limpieza, filtrado, proporción de mezcla, orden de entrenamiento y retroalimentación producen capacidades distintas en el modelo. La recipe es la barrera menos visible de la industria de datos. No consiste en «comprar un montón de datos», sino en decidir qué datos entran al entrenamiento, en qué proporción, en qué momento y cómo se evalúan.

\begin{figure}[H]
\centering
\includegraphics[width=0.92\textwidth]{data-recipe.png}
\caption{Data Recipe: fuentes, filtrado, proporción de mezcla, currículo y retroalimentación de la evaluación.}
\end{figure}

\begin{knowledgebox}{Lectura de la figura: por qué la recipe es más difícil de copiar que la materia prima}
En el centro de la figura está la Data Recipe y a su alrededor aparecen las fuentes, el filtrado, la proporción de mezcla, el currículo y la evaluación. La materia prima se puede comprar; la receta surge de la experimentación, los fracasos, la evaluación y el conocimiento del dominio. La recipe determina si los datos se convierten realmente en capacidad.
\end{knowledgebox}

\subsection{Términos clave: conceptos relacionados con la recipe}

\noindent\begin{tabularx}{\textwidth}{p{0.22\textwidth}p{0.32\textwidth}X}
\toprule
Término & Problema que resuelve & Significado en este episodio \\
\midrule
Filtering & Eliminar datos de baja calidad, duplicados, dañinos o irrelevantes & Determina la pureza de la señal de entrenamiento. \\
Mixture & Proporción entre datos de distintas fuentes & Determina la distribución de capacidades y las preferencias. \\
Curriculum & Orden en que los datos entran al entrenamiento & Determina qué aprende primero el modelo y qué aprende después. \\
Eval Feedback & Usar los resultados de la evaluación para deducir qué cambiar en los datos & Hace que la recipe sea iterable. \\
Synthetic Data & Datos generados por un modelo o por simulación & Compensa la escasez de datos reales, pero exige controlar el sesgo de distribución. \\
\bottomrule
\end{tabularx}

\subsection{Resumen de la sección}

La Data Recipe es el convertidor que transforma los datos en capacidad. Lo que de verdad resulta difícil de copiar en la industria de datos no suele ser la fuente de datos en sí, sino la receta y el ciclo cerrado de retroalimentación.
